\chapter*{Conclusion}
\addcontentsline{toc}{chapter}{Conclusion}

La régression logistique repose sur une chaîne cohérente : modèle linéaire des
log-cotes, loi de Bernoulli conditionnelle, maximum de vraisemblance, log-perte
convexe et optimisation numérique contrôlée. OvR réutilise exactement ce
problème binaire $K$ fois et fournit une règle multiclasses simple par maximum
des scores. Sa simplicité a une contrepartie précise : la somme des sorties
indépendantes est libre, alors qu'une distribution catégorielle exige une somme
égale à $1$. La validité du résultat dépend enfin autant du protocole
expérimental --- absence de fuite, test indépendant, diagnostics par classe ---
que de la correction des équations.

\begin{resultat}[Enchaînement des opérations]
Le modèle reçoit des caractéristiques préparées, calcule quatre combinaisons
linéaires, transforme éventuellement leurs scores par la fonction logistique et
retient le score maximal. Son apprentissage consiste à choisir chaque vecteur de
paramètres de façon à maximiser la vraisemblance de la cible binaire associée.
Sa validation consiste à vérifier séparément la convergence numérique, la
généralisation et les erreurs propres à chaque maison. Ces trois niveaux ne
doivent jamais être confondus.
\end{resultat}
